\documentclass{article}
\usepackage[T1]{fontenc}
\usepackage{lmodern}
\usepackage{graphicx}
\usepackage{amsmath,amssymb}
\usepackage[a4paper,margin=2cm]{geometry}
\usepackage[absolute,overlay]{textpos}
\setlength{\TPHorizModule}{1mm}
\setlength{\TPVertModule}{1mm}
\usepackage[indonesian]{babel}
\usepackage{hyperref}
\setlength{\emergencystretch}{2em}
% unit: Halaman 5

\begin{document}

% === Halaman 5 (python/native) ===

• Setiap blok plainteks dienkripsi dalam 16 putaran enciphering.

• Setiap putaran menggunakan kunci internal (kunci putaran) berbeda.

• Kunci internal sepanjang 48-bit dibangkitkan dari kunci eksternal

• Setiap blok plainteks mengalami permutasi awal (IP), 16 putaran 

enciphering, dan inversi permutasi awal (IP-1).  (lihat Gambar 1)

5

\end{document}
